% Part: sets-functions-relations
% Chapter: sets
% Section: basics

\documentclass[../../../include/open-logic-section]{subfiles}

\begin{document}

\olfileid[zh]{sfr}{set}{bas}
\olsection{外延性}

一个\emph{集合}是一批对象所构成的集合体，被视作一个单一
对象。构成集合的那些对象称为该集合的\emph{元素}或
\emph{成员}。若 $x$ 是集合~$A$ 的!!a{element}，则我们
记作 $x \in A$；若不然，则记作 $x \notin A$。不含任何!!{element}的集合称为\emph{空}集，记作~「」$\emptyset$''。

\begin{explain}
我们如何\emph{指定}该集合、如何\emph{排列}其!!{element}，乃至多少次\emph{计数}其!!{element}，都无关紧要。
要紧的只是其!!{element}到底是什么。我们将这一点落实到如下原则中。
\end{explain}

\begin{defn}[外延性]
  若 $A$ 与 $B$ 是集合，则 $A = B$ 当且仅当
  $A$ 的每个!!{element}也是 $B$ 的!!a{element}，反之
  亦然。
\end{defn}

外延性提供了若干记号。一般而言，当我们有一批
对象 $a_{1}$, \dots, $a_{n}$ 时，$\{a_{1}, \dots, a_{n}\}$ 便是
以 $a_1, \ldots, a_n$ 作为其!!{element}的\emph{那个}集合。我们强调
「\emph{那个}」这个词，因为外延性告诉我们只可能存在
\emph{一个}这样的集合。事实上，外延性还提供了下面这一条：
  \[
    \{a, a, b\} = \{a, b\} = \{b,a\}.
  \]
这也印证了如下观点：当我们考察集合时，我们并不关心其!!{element}的顺序，也不关心它们被指定了多少次。

\begin{tagblock}{novice}
\begin{ex}
每当你有一批对象时，你都可以把它们归拢到一个集合里。例如，理查德的兄弟姐妹所成的集合，就是一个只含一个人的集合，我们可以把它记作 $S=\{\textrm{露丝}\}$。小于 $4$ 的正整数所成的集合是 $\{1, 2, 3\}$，但也可以记作 $\{3, 2, 1\}$，甚至记作 $\{1, 2, 1, 2, 3\}$。依据外延性，这些都是同一个集合。因为 $\{1, 2, 3\}$ 的每个!!{element}也是 $\{3, 2, 1\}$（以及 $\{1,
2, 1, 2, 3\}$）的!!a{element}，反之亦然。
\end{ex}
\end{tagblock}

我们常常会借助其!!{element}所共有的某个性质来指定一个集合。为此，我们使用如下简写记号：
$\Setabs{x}{\phi(x)}$，其中 $\phi(x)$ 表示~$x$
为了被计入该集合的!!{element}之中所必须满足的性质。

\begin{tagblock}{novice}
\begin{ex}
在我们的例子里，我们也本可以把 $S$ 指定为
\[
S = \Setabs{x}{x \text{ 是理查德的兄弟姐妹}}.
\]
\end{ex}
\end{tagblock}

\begin{tagblock}{math}
\begin{ex}
一个数称为\emph{完全数}，当且仅当它等于其真因子（即能整除它、且与该数自身不相同的那些数）之和。例如，$6$ 是完全数，因为它的真因子是 $1$、$2$ 与~$3$，而 $6 = 1 + 2 + 3$。事实上，$6$ 是小于 $10$ 的唯一一个完全正整数。于是，利用外延性，我们可以说：
\[
	\{6\} = \Setabs{x}{x\text{ 是完全数且 }0 \leq x \leq 10}
\]
我们将右侧的记号读作「使得 $x$ 为完全数且 $0 \leq x \leq 10$ 的一切 $x$ 所成的集合」。这里的同一性印证了：当我们考察集合时，我们并不关心它们是如何被指定的。更一般地，外延性保证：使得 $\phi(x)$ 成立的 $x$ 所成的集合恰好只有一个。
因此，外延性使得我们可以把 $\Setabs{x}{\phi(x)}$ 称为使得~$\phi(x)$ 成立的 $x$ 所成的\emph{那个}集合。
\end{ex}
\end{tagblock}

外延性为我们提供了一种表明两个集合相同的方法：要表明 $A = B$，只需表明每当 $x \in A$ 时也有 $x \in B$，并且每当 $y \in B$ 时也有 $y \in A$。

\begin{prob}
证明至多只有一个空集，即证明：若 $A$ 与 $B$
都是不含任何!!{element}的集合，则 $A = B$。
\end{prob}

\end{document}
